\section{怎样写比较简单的议论文}

议论文是通过摆事实、讲道理来表达作者的见解和主张的一种文体。就初学写作的初中学生来说，应该先从比较简单的议论文练起。说它“比较简单”，是指论述的问题比较简单，论证的层次不多，论证的手法也比较单纯。它往往先确立一个明确的观点，然后用一两个事例或名言警句进行证明，从而得出结论。但是，它与记叙文相比较，就可看出它有自己的明显的特点：它主要运用逻辑推理直接阐明道理，反映客观事物的本质和规律，而不是主要通过形象地刻画人物、记叙事实来反映社会生活；它以议论为主要表达方式，而不是以叙述和描写为主要表达方式；它通常按“提出问题 --- 分析问题 --- 解决问题”的结构形式来议论说理，而不是通常按“开端 --- 发展 --- 高潮 --- 结局”的结构形式来记叙故事；它的语言要准确、概括，特别强调逻辑性，而不象记叙文那样作过多的描写和铺叙，较少形象的描绘。

我们在写比较简单的议论文时，应注意哪些问题呢？

1. 首先要确立一个正确、鲜明又有现实意义的论点。论点是文章的灵魂，议论的中心。一般说来，比较简单的议论文的题目往往就是论点。这类的题目经常用一个肯定语气的短语或单句形式出现，例如《为四化立志成材》、《锻炼贵在坚持》、《坚持长跑好》、《根深才能叶茂》、《失败是成功之母》等，提出的一种主张就是文章需要证明的论点。但是，有时题目只是提出一个问题，一个概念，或者提供一个事实，几句名言，例如《“开夜车”好不好？》、《小议“改过”》、《说“韧”》，或者提供象郭沫若接受演员的意见，把他的《屈原》剧本中的一句话“你是个无耻的文人”改成“你这个无耻的文人”的故事，而叫你根据这个故事写篇议论文，等等。这时就需要我们对有关的问题、论断或材料进行深入的分析，从中提炼、概括出自己写作的论点。例如该论关于郭沫若的这“一字之师”故事吧，我们可以提炼出几个方面的论点，象赞美长者的虚怀若谷，称道年轻人的敢想敢说，或者论述大师的成就来自严格的要求，甚至可以专讲锤炼语言的好处。在确立论点的过程中，我们不妨多提几个议题和论点，进行比较、选择，使最后确定的论点立意新一点，深一点，高一点。经过对事物的分析、概括而确立的论点，必须做到：要正确，符合马列主义基本原理，经得起社会实践的检验；要鲜明，态度明朗地对议论的问题表示赞成还是反对，不能模棱两可，含糊其辞；要有现实意义，提出人们普遍关心、迫切需要解决的问题来论述。

论点怎样表述呢？一般常用疑问句提出问题，针对这个问题用判断句的形式作出回答，这个回答就是论点；有时则干脆用一个判断句进行表述。例如初中语文课本第一册的两篇议论文，正是分别采用这两种方式进行表述的。毛泽东同志的《纪念白求恩》的论点是：“一个外国人，毫无利己的动机，把中国人民的解放事业当作他自己的事业，这是什么精神？这是国际主义的精神，这是共产主义的精神，每一个中国共产党员都要学习这种精神。”吴晗同志的《谈骨气》的论点是：“我们中国人是有骨气的。”应该注意：要使论点表述得鲜明突出，问题就要提得集中，切忌笼统，范围不宜太宽；作为整个议论的纲，又要能够概括议论的主要内容。论点的用语要准确严密，尽量简括。至于论点的位置则因文而异，不过初学写作者可以采用开宗明义就提出论点的方法，或在文章末尾明确地归纳出来，一般不要放在文章中间，这样有利于突出论点以及依据论点展开论述，可以防止忘写了论点或者脱离论点乱加发挥的毛病。

2. 围绕论点搜集和选择可靠、典型的论据。要让所说的道理使人理解和接受，必须要用适当的材料来说明它、证明它，这些材料就是论据。同学们爱用事例或史实（有时还用统计数字）来作论据，这当然是对的，因为事实胜于雄辩嘛。吴晗的《谈骨气》，就选用了文天祥宁为民族利益而死也不屈辱投降去做大官，古代一个穷人宁可饿死也不食“嗟来之食”，闻一多怒对国民党的手枪宁可倒下去也不愿屈服这样三个古今事例作论据，有力地证明了“我们中国人是有骨气的”这个论点。这种类型的论据叫做事实论据，它有很强的说服力，因为论点正确与否，归根结蒂要看其是否符合客观实际。有的同学还爱用为人们所公认的道理、名言警句或科学原理等来作论据，这当然也是对的，因为它们都是经过实践证实了的真理，也是符合客观实际的，同样具有很强的说服力。人们把这种类型的论据叫做理论论据。例如《谈骨气》一文中引用了战国时期思想家孟子的名言：“富贵不能淫，贫贱不能移，威武不能屈，此之谓大丈夫。”不正有力地说明了我们中国人自古以来就是有骨气的吗？

应该注意，选用的论据必须做到可靠而又典型。可靠是指材料要真实：事实确凿，经得起核对，每个细节、每个数字都不允许为了附会观点而掺进浮夸不实的成分；引用完整、准确，忠于原文意思，不能断章取义，更不允许歪曲篡改，胡乱编造。典型是指材料要有代表性，能反映事物的本质，因为有些事情虽然属实，但它只是生活中偶然的个别的现象，不能反映事物的本质，罗列再多也不能使论点具有说服力。例如《谈骨气》引用的孟子的话一字不错，列举的关于文天祥、饿者和闻一多的三个事例有史可据，其中决没有想当然的“可能”、“也许”或“有的”、“某些”之类含糊不清的字眼，更没有张冠李戴、违反常识的错误，这就叫可靠，富有说服力；同时，所用的三个事例并不是随便想到就用的，而是经过精心的选择，象文天祥的事例正是“富贵不能淫”的典型，饿者的事例正是“贫贱不能移”的典型，闻一多的事例正是“威武不能屈”的典型，这样典型的事例能反映事物的本质，说服力极强，使文章的论点不可动摇。

另外，作为议论文的事实论据，表述时一定要抓住最能证明论点的关键部分作简明扼要的概述，而不要用记叙的笔调作具体的叙述，更不能写成有头有尾的故事，因为议论文中用到的事实，是作为论据出现的。如果这事实是人所共知的，那只要明确地点一下，使它能为证明自己的论点服务即可，即使这事实人们比较生疏，需要多写几句才能使读者明白，也应抓住要点，概括表达，否则，就会以叙代议，喧宾夺主。象《谈骨气》所举的三个古今人物的典型事例，都只抓住他们在生死关头、宁死不屈的最能说明论点的那一段事迹写，文字十分概括，简要，至于他们平生中其他动人的事迹则一律舍弃了。这样的写法，我们应该好好体会，认真学习。

3.论证要掌握基本方法，使论据有力地证明论点的正确性。大家知道，论点解决“要证明什么”的问题，论据解决“用什么证明”的问题，论证则解决“如何进行证明”的问题，即如何用论据来说明论点，揭示论点和论据之间的逻辑关系，从而把问题分析透彻。议论文从论证的方式看，一般可以分为立论和驳论两种。写比较简单的议论文，首先就得掌握立论的基本方法。

所谓立论，是就一定的事件或问题，从正面提出并阐明自己的论点。立论的方法主要有例证、引证（引用名言警句、科学公理、人所共知的常理等作论据，以证明论点的正确性）、比较（用过去的情况跟当前的情况作纵向的比较，或者用两种对立的事物作横的比较，从比较中显示事物的本质，证明论点的正确性）、类比（从已知的个别事例推论到具有某些相同性质的个别事例，往往通过讲个故事、举个类似的例子，打个比方或借助一些成语典故等，来说明抽象的道理，证明自己的论点），分析等。其中以例证和分析为两种最常见、最基本的论证方法，而例证尤为重要。

例证法，是一种用事例作为论据来证明论点正确性的论证方法。这种方法，大家一定非常熟悉，在现实生活中也经常运用。例如一位同学就谁可当三好学生发言说：“我赞成把××同学评为三好学生。他积极要求进步，努力搞好他所担任的卫生委员工作，已被团委批准加入共青团，他不仅自己认真刻苦地学好各门功课，还能主动地帮助同学；在体育锻炼中，他还是个积极分子，已经全部达到国家三级劳动锻炼标准。这样的同学怎能不评为三好学生呢？”这个发言，实际是一篇很短的简单议论文：第一句先提出自己的主张（论点）；接着用三方面的事实说明自己的理由（论据）；最后一句以进一步强调论点作结。他的论证方法，就是提出论点，运用可靠、典型的事例进行证明，最后得出结论。在结构上是按照“总 --- 分 --- 总”的方式组织的。这样的论证方法就是例证法。通过例证，使论点和论据统一起来，结合成为一个整体，达到说服读者的目的。如果根据这个发言，把论据的事实充实起来，使之具体化，再加上必要的分析，就能写成一篇比较简单的议论文。

应该注意：在运用事实论据证明论点时，一定要使论据与论点之间具有严密的逻辑联系。如果论据与论点对不上口径，或者论据只能证明论点的某一部分，那么论据罗列再多，也是不能证明论点的。例如，假若那位同学所提的理由（论据），只是说××同学热心搞家务劳动，积极参加街道业余剧团，是个谁都打不过的摔跤能手，平时虽不抓紧学习，但是考试成绩总还不错，这一切与三好学生的标准并不吻合，能说服得了全班同学吗？当然不能。即使那位同学提出××同学怎样做好班务工作，怎样积极参加校内外社会活动和义务劳动，怎样爱护公物，怎样助人为乐，却不提怎样搞学习，怎样积极锻炼身体，结果“三好”只有“一好”，这样的论据当然就不能全面地论证论点的正确性，因此也缺乏说服力。

我们再谈谈分析法。分析法是一种用分析问题、剖析事理来揭示论点和论据之间的因果关系，证明论点正确性的方法。我们知道，为了使论据有力地证明论点的正确性，就得摆事实、讲道理。例证法是摆事实的主要手段，而分析法就是讲道理的主要手段。常常听到有人说议论文不好写，其中很重要的一个原因就是不善于分析。比如叫写一篇《要珍惜时间》的议论文，有的同学觉得没啥可写的。他也想，珍惜时间，就能搞好学习；不珍惜时间，就不能搞好学习。把这个想法写下来，文章显得干瘪而空洞，即使加上珍惜时间而搞好学习的事例，文章也还显得单薄而肤浅。善于分析的同学就不同了，他可以从青年学生为什么要珍惜时间和怎样珍惜两个方面去展开议论。这两个方面又可以分析，如为什么要珍惜时间呢？一是因为时间就是生命；二是时间稍纵即逝，一去不返；三是学生时代尤其短暂，尤其可贵，等等。又如怎样珍惜时间呢？一是学习要有计划，时间安排合理；二是学习中要抓住点滴时间；三是要改进学习方法，提高时间的利用率，等等。这样一分析，说理就容易透彻，内容也就显得充实起来了。

从这个小例子可以看出，开始分析，需要从提出问题入手，推进分析，使分析进一步深入，同样也需要提出问题。一般地说，提问有两种类型：一是从事物的各个方面（即从横的方面）去提出问题，如“是什么”、“怎么样”、“为什么”；二是从事物的某一方面（即从纵的方面）去刨根究底。要使问题提得恰当，分析中肯，就得以事实为基础，以自己的亲身观察和体验为依据。毛泽东同志说：“所谓分析，就是分析事物的矛盾。不熟悉生活，对于所论的问题不真正了解，就不可能有中肯的分析。”因此，我们在分析时，需要占有事实材料，根据事实来研究与认识事物的实质，而不要空发议论。这就是说，开始学写议论文，最好把分析和例证结合起来运用，使事实摆得充分、典型，使道理讲得明白、深透。光有例证而无分析，不容易揭示论点和论据之间的因果关系，说理不透；光有分析而无例证，也不容易使论点很好地确立起来，使议论具有巨大的说服力。
\newpage